\chapter{El Incidente}
\label{chap:incidente}

\section{Mapeo del Ataque a la Cyber Kill Chain y MITRE ATT\&CK}

En esta sección se presenta una referencia técnica que describe el ataque desde la perspectiva del adversario. Se detallan las acciones realizadas por el atacante, los objetivos perseguidos en cada fase y las técnicas MITRE ATT\&CK empleadas, organizadas conceptualmente de forma que **cada tabla corresponde a una Táctica de MITRE ATT\&CK** dentro de su respectiva fase de la Cyber Kill Chain. Además, cada acción se relaciona con los artefactos del sistema víctima donde pueden encontrarse evidencias relevantes, proporcionando una guía práctica para el análisis forense y la reconstrucción del incidente.

La Cyber Kill Chain se utiliza como marco para organizar la exposición de las distintas fases del ataque, manteniendo el orden cronológico natural en el que fueron ejecutadas. Sin embargo, es importante destacar que este orden no suele coincidir con el flujo real de una investigación \acrshort{dfir}. En la práctica, el analista raramente dispone de visibilidad sobre las fases iniciales de la intrusión; por el contrario, la investigación suele comenzar a partir de los indicadores de impacto o de las acciones finales del atacante. A partir de estas evidencias iniciales, el analista va reconstruyendo progresivamente la cadena de eventos, siguiendo el rastro de artefactos y evidencias hacia etapas anteriores del compromiso.

Por este motivo, aunque las fases se presentan siguiendo la secuencia clásica de la Cyber Kill Chain, una lectura en sentido inverso refleja de forma más fiel el proceso de investigación que seguiría un analista para reconstruir la historia completa del ataque.

El objetivo de esta sección es ofrecer una guía técnica estructurada que relacione de forma directa:
\begin{itemize}
    \item La fase de la Cyber Kill Chain,
    \item La Táctica MITRE ATT\&CK asociada a cada objetivo táctico,
    \item Las Técnicas concretas desarrolladas dentro de dicha táctica y su procedimiento técnico,
    \item Los artefactos del sistema donde podría hallarse evidencia forense de cada acción.
\end{itemize}

\subsection{Resumen de Fases del Ataque}

El desglose general y resumen de las distintas fases que componen la intrusión se encuentra detallado en el capítulo de Metodología (véase la Sección~\ref{sec:indice-fases-ataque}). 

A continuación, en las siguientes secciones de este capítulo, se analiza en detalle cada una de dichas fases, profundizando en las acciones específicas del atacante organizadas por Tácticas de MITRE ATT\&CK y los artefactos forenses concretos donde se manifiestan sus evidencias en el sistema comprometido.




\subsection{Reconocimiento} \label{sec:reconocimiento}

En esta fase, el atacante intenta recopilar información relevante sobre la organización, sus sistemas y posibles vectores de entrada. Desde el punto de vista del analista, estas actividades suelen considerarse parte de un análisis preventivo más que de un análisis de incidente propiamente dicho.

Dentro de esta categoría se incluyen tareas como el monitoreo de actividades de reconocimiento dirigidas a los endpoints expuestos de la organización, la revisión de bases de datos comercializadas por ciberdelincuentes que contengan información relacionada con la empresa y, en general, la supervisión de cualquier actividad externa que pueda indicar la recopilación de información previa a un ataque.



\subsection{Armamento} \label{sec:armamento}

En esta fase, el atacante prepara los \textit{payloads}, herramientas y la infraestructura que utilizará durante el ataque. Al igual que ocurre con la fase de reconocimiento, esta etapa se sitúa más cerca del análisis preventivo que del análisis forense de un incidente.

Las acciones llevadas a cabo por el atacante en esta fase se describen en el capítulo \textit{Red Team Notes} (ver Sección~\ref{chap:red-team}), donde se detallan, desde la perspectiva ofensiva, el desarrollo de los \textit{payloads}, las decisiones técnicas tomadas y la infraestructura empleada para ejecutar el ataque.

Desde el punto de vista del analista, esta fase puede apoyarse en servicios externos que realizan escaneos continuos de Internet en busca de infraestructuras maliciosas, como servidores de \textit{Cobalt Strike} u otros \textit{\acrshort{c2}} con configuraciones mínimas o pobre \textit{fingerprinting}. Estos servicios permiten a la organización detectar tempranamente posibles activos maliciosos relacionados con campañas activas y, en caso necesario, aplicar medidas preventivas como el bloqueo automático de conexiones hacia o desde dichas infraestructuras.


\newpage
\subsection{Entrega} \label{sec:entrega}

En la fase de Entrega, el atacante interactúa por primera vez con la organización víctima marcando el inicio del incidente. El objetivo es introducir el primer \textit{payload} en el entorno objetivo mediante un dispositivo físico \textit{BadUSB}.

Para un análisis detallado de las decisiones ofensivas tomadas durante esta fase, véase el Capítulo~\ref{chap:red-team}, Sección~\ref{sec:bad-usb}.

\subsubsection{Mapeo por Tácticas MITRE ATT\&CK}

A continuación se presentan las tácticas desarrolladas en esta fase. En la Tabla~\ref{tab:tactica-initial-access-entrega} se aborda la táctica de Acceso Inicial, mientras que en la Tabla~\ref{tab:tactica-execution-entrega} se detalla la táctica de Ejecución asociada al dispositivo.

\begin{table}[htbp]
    \centering
    \caption{Táctica: Initial Access (TA0001) — Fase de Entrega}
    \label{tab:tactica-initial-access-entrega}
    \renewcommand{\arraystretch}{1.3}
    \small
    \begin{tabularx}{\textwidth}{|>{\bfseries\columncolor{gray!15}}p{3.8cm}|X|}
        \hline
        \multicolumn{2}{|c|}{\cellcolor{blue!15}\bfseries Táctica: Initial Access (TA0001)} \\ \hline
        \textbf{Técnica} & \textbf{Procedimiento y Ejemplo Concreto} \\ \hline
        Replication Through Removable Media (T1091) & El atacante conecta físicamente un dispositivo \textit{BadUSB} programado en el equipo víctima. El dispositivo es reconocido por el sistema operativo como una unidad removible y periférico de entrada legítimo, superando los controles de bloqueo de almacenamiento USB convencionales. \\ \hline
    \end{tabularx}
\end{table}

\begin{table}[htbp]
    \centering
    \caption{Táctica: Execution (TA0002) — Fase de Entrega}
    \label{tab:tactica-execution-entrega}
    \renewcommand{\arraystretch}{1.3}
    \small
    \begin{tabularx}{\textwidth}{|>{\bfseries\columncolor{gray!15}}p{3.8cm}|X|}
        \hline
        \multicolumn{2}{|c|}{\cellcolor{blue!15}\bfseries Táctica: Execution (TA0002)} \\ \hline
        \textbf{Técnica} & \textbf{Procedimiento y Ejemplo Concreto} \\ \hline
        Input Injection (T1674) & Tras ser reconocido, el \textit{BadUSB} emula un teclado (\acrshort{hid}) e inyecta de forma automatizada una secuencia de pulsaciones de teclas: ejecuta la combinación \texttt{Win + R}, escribe \texttt{powershell.exe} y lanza la ejecución de comandos para descargar y ejecutar el script inicial sin interacción del usuario. \\ \hline
    \end{tabularx}
\end{table}

\subsubsection{Artefactos relevantes}

\begin{itemize}
    \item \textbf{Event Log (System / DriverFrameworks-UserMode)}: Registra la conexión y enumeración del dispositivo USB, incluyendo el id de dispositivo y la marca temporal de conexión, permitiendo demostrar la presencia física del BadUSB.
    \item \textbf{Sysmon (EID~1 — Process Create)}: Registra la creación del proceso \texttt{powershell.exe} invocada desde \texttt{explorer.exe} tras la combinación de teclas inyectada.
\end{itemize}


\newpage
\subsection{Explotación} \label{sec:explotacion}

En la fase de Explotación (Stage 1), el adversario consigue la primera ejecución de código malicioso e inicia el establecimiento de una conexión remota interactiva (\textit{reverse shell}) con el servidor atacante.

Para una descripción detallada del desarrollo y funcionamiento de este \textit{payload}, véase el Capítulo~\ref{chap:red-team}, Sección~\ref{sec:shellcode}.

\subsubsection{Mapeo por Tácticas MITRE ATT\&CK}

Esta fase abarca dos tácticas principales: Ejecución (Tabla~\ref{tab:tactica-execution-explotacion}) y Mando y Control (Tabla~\ref{tab:tactica-c2-explotacion}).

\begin{table}[htbp]
    \centering
    \caption{Táctica: Execution (TA0002) — Fase de Explotación}
    \label{tab:tactica-execution-explotacion}
    \renewcommand{\arraystretch}{1.3}
    \small
    \begin{tabularx}{\textwidth}{|>{\bfseries\columncolor{gray!15}}p{4.2cm}|X|}
        \hline
        \multicolumn{2}{|c|}{\cellcolor{blue!15}\bfseries Táctica: Execution (TA0002)} \\ \hline
        \textbf{Técnica} & \textbf{Procedimiento y Ejemplo Concreto} \\ \hline
        Command and Scripting Interpreter: PowerShell (T1059.001) & Ejecución del script desofuscado en PowerShell que prepara el entorno, evade AMSI en memoria e invoca las instrucciones de descarga del payload secundario. \\ \hline
        User Execution: Malicious File (T1204.002) & Ejecución encubierta del binario descargado mediante el parámetro \texttt{-WindowStyle Hidden}, iniciando la payload en segundo plano. \\ \hline
    \end{tabularx}
\end{table}

\begin{table}[htbp]
    \centering
    \caption{Táctica: Command and Control (TA0011) — Fase de Explotación}
    \label{tab:tactica-c2-explotacion}
    \renewcommand{\arraystretch}{1.3}
    \small
    \begin{tabularx}{\textwidth}{|>{\bfseries\columncolor{gray!15}}p{4.2cm}|X|}
        \hline
        \multicolumn{2}{|c|}{\cellcolor{blue!15}\bfseries Táctica: Command and Control (TA0011)} \\ \hline
        \textbf{Técnica} & \textbf{Procedimiento y Ejemplo Concreto} \\ \hline
        Ingress Tool Transfer (T1105) & El script descarga un archivo temporal (\texttt{2fd19d11-d4b7-46d1-94db-7edd16980d15.tmp}) desde el servidor del atacante hacia la ruta \texttt{\$env:TEMP} y lo renombra a \texttt{.exe}. \\ \hline
        Remote Services / Reverse Shell (T1021) & El binario ejecutado abre un socket TCP saliente hacia el puerto de escucha del atacante, proporcionando una shell de comandos interactiva. \\ \hline
    \end{tabularx}
\end{table}

\subsubsection{Artefactos relevantes}

\begin{itemize}
    \item \textbf{PowerShell Operational Log (EID~4104 — ScriptBlockLogging)}: Registra el contenido íntegro del script ejecutado, incluyendo la URL de descarga y los comandos de ejecución.
    \item \textbf{Sysmon (EID~1 — ProcessCreate / EID~3 — NetworkConnect)}: El EID~1 captura la ejecución del binario en \texttt{\$env:TEMP} y el EID~3 registra la conexión saliente de la reverse shell.
    \item \textbf{\$MFT}: Registra las marcas temporales MACB de creación y renombrado del fichero ejecutable en el directorio temporal.
\end{itemize}


\newpage
\subsection{Instalación} \label{sec:instalacion}

En la fase de Instalación (Stage 2), el atacante consolida su presencia en la máquina mediante técnicas de \textit{\acrshort{dll} Hijacking}, \textit{Side Loading} y \textit{\acrshort{dll} Proxying}, asegurando asimismo la persistencia en el sistema.

Un análisis en profundidad de estos procedimientos ofensivos se encuentra en el Capítulo~\ref{chap:red-team}, Secciones~\ref{sec:dll-hijacking} y~\ref{sec:persistencia}.

\subsubsection{Mapeo por Tácticas MITRE ATT\&CK}

Esta fase comprende las tácticas de Evasión de Defensas (Tabla~\ref{tab:tactica-defense-evasion-instalacion}), Ejecución (Tabla~\ref{tab:tactica-execution-instalacion}) y Persistencia (Tabla~\ref{tab:tactica-persistence-instalacion}).

\begin{table}[htbp]
    \centering
    \caption{Táctica: Defense Evasion (TA0005) — Fase de Instalación}
    \label{tab:tactica-defense-evasion-instalacion}
    \renewcommand{\arraystretch}{1.3}
    \small
    \begin{tabularx}{\textwidth}{|>{\bfseries\columncolor{gray!15}}p{4.2cm}|X|}
        \hline
        \multicolumn{2}{|c|}{\cellcolor{blue!15}\bfseries Táctica: Defense Evasion (TA0005)} \\ \hline
        \textbf{Técnica} & \textbf{Procedimiento y Ejemplo Concreto} \\ \hline
        File and Directory Permissions Modification (T1222) & Copia del ejecutable legítimo de Calibre a \texttt{\%LOCALAPPDATA\%\textbackslash Microsoft\textbackslash WindowsApps}, un directorio con permisos de escritura para el usuario donde se controla el orden de carga de librerías. \\ \hline
        Hijack Execution Flow: DLL Search Order Hijacking (T1574.001) & Ubicación de la \acrshort{dll} maliciosa \texttt{calibre-launcher.dll} en el mismo directorio del ejecutable para forzar su carga con preferencia sobre la librería de sistema. \\ \hline
        Masquerading: Match Legitimate Name or Location (T1036.005) & Renombrado de la \acrshort{dll} legítima a \texttt{calibre-launcher-old.dll} y suplantación de su nombre original por la librería maliciosa que actúa como proxy reenviando las llamadas legítimas. \\ \hline
    \end{tabularx}
\end{table}

\begin{table}[htbp]
    \centering
    \caption{Táctica: Execution (TA0002) — Fase de Instalación}
    \label{tab:tactica-execution-instalacion}
    \renewcommand{\arraystretch}{1.3}
    \small
    \begin{tabularx}{\textwidth}{|>{\bfseries\columncolor{gray!15}}p{4.2cm}|X|}
        \hline
        \multicolumn{2}{|c|}{\cellcolor{blue!15}\bfseries Táctica: Execution (TA0002)} \\ \hline
        \textbf{Técnica} & \textbf{Procedimiento y Ejemplo Concreto} \\ \hline
        Ingress Tool Transfer (T1105) & Transferencia y descarga desde el \acrshort{c2} de la \acrshort{dll} maliciosa \texttt{calibre-launcher.dll} al directorio de trabajo del binario de Calibre. \\ \hline
    \end{tabularx}
\end{table}

\begin{table}[htbp]
    \centering
    \caption{Táctica: Persistence (TA0003) — Fase de Instalación}
    \label{tab:tactica-persistence-instalacion}
    \renewcommand{\arraystretch}{1.3}
    \small
    \begin{tabularx}{\textwidth}{|>{\bfseries\columncolor{gray!15}}p{4.2cm}|X|}
        \hline
        \multicolumn{2}{|c|}{\cellcolor{blue!15}\bfseries Táctica: Persistence (TA0003)} \\ \hline
        \textbf{Técnica} & \textbf{Procedimiento y Ejemplo Concreto} \\ \hline
        Scheduled Task/Job (T1053) & Creación de la tarea programada \texttt{calibre-update} configurada para ejecutarse al inicio del sistema. La tarea invoca \texttt{calibre-debug.exe}, provocando la carga automática de la \acrshort{dll} maliciosa tras cada reinicio. \\ \hline
    \end{tabularx}
\end{table}

\subsubsection{Artefactos relevantes}

\begin{itemize}
    \item \textbf{\$MFT}: Registro de las operaciones de copia, renombrado de DLLs y timestamps de creación de ficheros en el directorio del usuario.
    \item \textbf{Sysmon (EID~11 — FileCreate / EID~7 — ImageLoad)}: Registra la creación de la DLL maliciosa y su posterior carga por parte del proceso legítimo de Calibre.
    \item \textbf{TaskScheduler Operational Log (EID~106, 200, 201)}: Documenta el registro y ejecución de la tarea programada \texttt{calibre-update}.
\end{itemize}


\newpage
\subsection{Mando y Control (C2)} \label{sec:comando-control}

En la fase de Comando y Control, el adversario mantiene una sesión activa y persistente en memoria mediante el implante de \textit{Sliver C2}, utilizando comunicaciones cifradas y técnicas de inyección en memoria.

Para una descripción detallada sobre la carga reflectiva en memoria y la infraestructura C2, consúltese el Capítulo~\ref{chap:red-team}, Secciones~\ref{sec:reflective} y~\ref{sec:infraestructura}.

\subsubsection{Mapeo por Tácticas MITRE ATT\&CK}

Esta fase se estructura en las tácticas de Evasión de Defensas (Tabla~\ref{tab:tactica-defense-evasion-c2}) y Mando y Control (Tabla~\ref{tab:tactica-c2-c2}).

\begin{table}[htbp]
    \centering
    \caption{Táctica: Defense Evasion (TA0005) — Fase de Comando y Control}
    \label{tab:tactica-defense-evasion-c2}
    \renewcommand{\arraystretch}{1.3}
    \small
    \begin{tabularx}{\textwidth}{|>{\bfseries\columncolor{gray!15}}p{4.2cm}|X|}
        \hline
        \multicolumn{2}{|c|}{\cellcolor{blue!15}\bfseries Táctica: Defense Evasion (TA0005)} \\ \hline
        \textbf{Técnica} & \textbf{Procedimiento y Ejemplo Concreto} \\ \hline
        Reflective Code Loading (T1620) & La \acrshort{dll} maliciosa descarga e inyecta el implante Sliver directamente en memoria mediante carga reflectiva (\textit{unbacked memory}), sin tocar disco ni utilizar las APIs convencionales de carga de módulos (\texttt{LoadLibrary}). \\ \hline
    \end{tabularx}
\end{table}

\begin{table}[htbp]
    \centering
    \caption{Táctica: Command and Control (TA0011) — Fase de Comando y Control}
    \label{tab:tactica-c2-c2}
    \renewcommand{\arraystretch}{1.3}
    \small
    \begin{tabularx}{\textwidth}{|>{\bfseries\columncolor{gray!15}}p{4.2cm}|X|}
        \hline
        \multicolumn{2}{|c|}{\cellcolor{blue!15}\bfseries Táctica: Command and Control (TA0011)} \\ \hline
        \textbf{Técnica} & \textbf{Procedimiento y Ejemplo Concreto} \\ \hline
        Application Layer Protocol: Web Protocols (T1071.001) & El implante Sliver mantiene la comunicación con el servidor C2 mediante peticiones HTTPS periódicas (\textit{beaconing}) simulando tráfico de navegación web legítimo. \\ \hline
        Encrypted Channel (T1573) & Comunicación cifrada de extremo a extremo que impide la inspección directa del contenido de los comandos ejecutados por el atacante. \\ \hline
    \end{tabularx}
\end{table}

\subsubsection{Artefactos relevantes}

\begin{itemize}
    \item \textbf{Volcado de Memoria RAM}: Identificación del implante de Sliver en regiones de memoria ejecutable sin respaldo en disco (\textit{unbacked memory}).
    \item \textbf{Sysmon (EID~3 — NetworkConnect)}: Conexión de red saliente iniciada desde el proceso anfitrión hacia la IP del servidor C2.
    \item \textbf{Suricata Log}: Detección de patrones de comunicación C2, certificados TLS autofirmados y firmas JA3/JA3S anómalas.
\end{itemize}


\newpage
\subsection{Acciones sobre el objetivo} \label{sec:acciones-sobre-objetivo}

En la fase final de la intrusión, el atacante ejecuta una acción controlada para simular el impacto en el equipo comprometido sin provocar daños irreversibles.

\subsubsection{Mapeo por Tácticas MITRE ATT\&CK}

En la Tabla~\ref{tab:tactica-impact-acciones} se especifica la táctica de Impacto correspondiente a esta fase final.

\begin{table}[htbp]
    \centering
    \caption{Táctica: Impact (TA0040) — Fase de Acciones sobre el Objetivo}
    \label{tab:tactica-impact-acciones}
    \renewcommand{\arraystretch}{1.3}
    \small
    \begin{tabularx}{\textwidth}{|>{\bfseries\columncolor{gray!15}}p{4.2cm}|X|}
        \hline
        \multicolumn{2}{|c|}{\cellcolor{blue!15}\bfseries Táctica: Impact (TA0040)} \\ \hline
        \textbf{Técnica} & \textbf{Procedimiento y Ejemplo Concreto} \\ \hline
        Data Manipulation (T1565) & Despliegue e invocación en memoria de un assembly .NET desde el implante de Sliver que genera el archivo \texttt{desktop\_proof.txt} en el escritorio del usuario como indicador de compromiso (IOC), simulando el impacto de un ataque de ransomware sin destruición de datos. \\ \hline
    \end{tabularx}
\end{table}

\subsubsection{Artefactos relevantes}

\begin{itemize}
    \item \textbf{Sysmon (EID~11 — FileCreate / EID~7 — ImageLoad)}: Captura de la creación del fichero \texttt{desktop\_proof.txt} en el escritorio y carga de la infraestructura del .NET CLR en el proceso anfitrión.
    \item \textbf{.NET CLR Logs (ETW)}: Evidencia de la carga dinámica del assembly inyectado en memoria.
    \item \textbf{Volcado de Memoria RAM}: Extracción del assembly .NET inyectado para su análisis e inspección de código.
\end{itemize}
